\begin{theorem}[Directed transport of reconstructed internal bonds]
    \label{thm:peps_internal_gauge_transport}
    \lean{TNLean.PEPS.regularDirectedTransport_eq_of_internal_reconstruction,TNLean.PEPS.regularDirectedTransport_injective_on_edge,TNLean.PEPS.regularDirectedTransport_eq_of_endpoints,TNLean.PEPS.regularRegionBondExtension_congr_of_internal,TNLean.PEPS.regularRegionBondExtension_eq_of_internal}
    \leanok
    Let $\Gamma$ be a simple graph with ordered vertices, $R$ a finite vertex
    set, and $k_v\in G$ for $v\in R$. If its internal ordered bond assignments
    satisfy $w_e=k_{\mathrm{head}(e)}u_e k_{\mathrm{tail}(e)}^{-1}$, then
    the directed transport along any internal step $x\to y$ satisfies
    \begin{align}
        D_w(x\to y)=k_yD_u(x\to y)k_x^{-1}.
    \end{align}
    Directed transport determines its ordered bond coefficient and depends
    only on the endpoints of the step. Equal internal assignments give equal
    extensions with any common exterior assignment; in particular, extending
    an internal restriction of an ambient assignment returns that assignment.
    These are auxiliary identities for the accessible-coordinate construction
    in \cite[Section~6.3]{Schuch2010PEPS}.
\end{theorem}
\begin{proof}
    \leanok
    \uses{thm:peps_regular_walk_holonomy}
    A reversed traversal uses the inverse operator, which gives the stated
    formula in both possible endpoint orders. The remaining assertions follow
    by separating internal and exterior bonds.
\end{proof}
